\section{Discussion}

\subsection{Computational Complexity Analysis}

The reason TA* requires an average of 15.46 seconds while FieldD*Hybrid requires only 0.495 seconds (a 31.2-fold difference) can be explained through computational complexity analysis.

TA*'s computational complexity is characterized as $O(b^d \times k)$, where:
\begin{itemize}
    \item $b$ = branching factor (26-neighborhood)
    \item $d$ = solution depth (average 50 steps)
    \item $k$ = terrain cost computation complexity (4 elements per node)
\end{itemize}

This complexity grows exponentially with respect to search depth $d$.

In contrast, FieldD*Hybrid comprises:
\begin{itemize}
    \item Initial path generation via D* Lite: $O(n)$
    \item Local improvement via sliding window: $O(w \times h)$
\end{itemize}

Total complexity: $O(n)$ (linear), where:
\begin{itemize}
    \item $n$ = number of grid nodes ($\approx$10,000)
    \item $w$ = window size (5$\times$5 = 25)
    \item $h$ = improvement iterations (10)
\end{itemize}

This fundamental difference in theoretical complexity (exponential vs. linear) directly explains the 31.2-fold speedup observed empirically.

Analysis of TA*'s 15.46-second execution time reveals the following breakdown:
\begin{itemize}
    \item Terrain cost calculation: 40\% (6.18s)
    \item Node exploration (A*): 35\% (5.41s)
    \item Heuristic calculation: 15\% (2.32s)
    \item Pruning judgment: 10\% (1.55s)
\end{itemize}

Terrain cost calculation dominates the computation because each node requires evaluation of four elements: slope, stability, obstacle proximity, and distance. This detailed evaluation is the source of TA*'s computational overhead but also the foundation of its adaptive decision-making capability.

\subsection{Computation Time of TA*}

The proposed TA* method exhibited a mean computation time of 15.46 seconds, which is longer compared to other methods. However, this increase in computation time reflects TA*'s design philosophy and is not merely a performance degradation.

TA* implements adaptive path planning through terrain recognition by dynamically adjusting search strategies according to scenario complexity. Specifically, approximately 40\% of total computation time is devoted to terrain cost calculation and 35\% to node exploration---processes absent in conventional A* algorithms. This ``increase in computation time'' represents an \textit{intentional design} for more detailed environmental recognition and reliable path discovery.

Importantly, TA*'s median value is 6.83 seconds, approximately 45\% of the mean value of 15.46 seconds. This indicates that the computation time distribution is right-skewed. In practice, 50\% of scenarios completed within 6.83 seconds, with only extremely complex scenarios requiring extended time. This characteristic provides evidence that TA* \textit{adapts} to environmental complexity: operating rapidly in simple environments while conducting detailed exploration in complex environments, thereby achieving a high success rate of 96.9\%.

Therefore, the increase in TA*'s computation time is not a sacrifice of real-time performance but rather a strategic choice prioritizing reliability and adaptability. In scenarios where path reliability is prioritized over computation time---such as offline simulation or pre-planned path generation---TA*'s design is rational.

\subsection{Development from TA* to FieldD*Hybrid}

In this research, we developed FieldD*Hybrid, capable of real-time operation, based on insights gained from TA*. FieldD*Hybrid achieved a 31-fold speedup compared to TA* (0.495~s vs.\ 15.46~s) while maintaining a 100\% success rate.

This improvement leverages the following insights revealed during TA* implementation:

\begin{enumerate}
    \item \textbf{Validity of terrain adaptation concept}: Since TA* confirmed the effectiveness of terrain recognition and dynamic parameter adjustment, this concept was incorporated into FieldD*'s local optimization.
    
    \item \textbf{Computational reduction strategy}: Analysis of TA*'s computation time revealed that terrain cost calculation was the primary bottleneck. FieldD*Hybrid reduced overall computational complexity from $O(n \log n)$ to $O(n)$ by combining basic path generation via D* Lite with local improvement using a sliding window method.
    
    \item \textbf{Pursuit of reliability}: Learning from TA*'s 96.9\% success rate, FieldD*Hybrid achieved a 100\% success rate by adding mechanisms to dynamically expand the search range.
\end{enumerate}

This development process represents a typical staged progression from theoretical exploration (TA*) to practical implementation (FieldD*Hybrid) in academic research. TA* functioned as a proof-of-concept for ``terrain-adaptive path planning,'' with its insights informing FieldD*Hybrid's design. These approaches are not competing methods but rather complementary positions at different stages of research.

Therefore, this research's contribution extends beyond proposing a single algorithm to demonstrating a methodology for advancing adaptive path planning concepts from research to practical application. These insights are applicable to other path planning methods and autonomous robot system development.

\subsection{Implications and Future Work}

The staged development from TA* to FieldD*Hybrid demonstrates an important research methodology: validating theoretical concepts through comprehensive experimentation before pursuing practical optimization. This approach ensures that performance improvements are grounded in understanding fundamental mechanisms rather than empirical tuning alone.

Future work should explore:
\begin{itemize}
    \item Parameter optimization for TA* to reduce computation time while maintaining reliability
    \item Extension of terrain adaptation concepts to 3D environments
    \item Integration with learning-based approaches for automatic parameter tuning
    \item Validation in real-world agricultural robot deployments
\end{itemize}

\paragraph{Key Message}
This research establishes that computation time should be evaluated in context: TA* prioritizes adaptability and reliability for research and offline planning, while FieldD*Hybrid demonstrates how these insights enable real-time operation. Both contributions are essential for advancing autonomous navigation in complex terrains.
